\section{Validating the Loser-Side Ranking}
\label{sec:validation}

This section validates the loser-side ranking, not the legal
status of ranked firms. The target is loser-side cartel
adjacency: always-loser firms observed alongside direct CADE
defendants in adjudicated procurement-cartel environments. That
target follows from the data layer. Award records observe
participation and outcomes; they do not observe agreements,
communications, or tender-specific conduct. The question is
therefore not whether the screen proves cartel membership, but
whether it concentrates the loser-side firms most plausibly
worth bid-layer forensic attention.

We organize the evidence as a validation architecture, with one
exercise per threat to interpretation. The conservative
adjudication benchmark
(subsection~\ref{sec:validation_conservative}) anchors the
target in cases resolved before the end of 2020, reducing
reliance on later case outcomes. The temporal holdout
(subsection~\ref{sec:validation_temporal}) asks whether the
ranking discriminates later-adjudicated cobidders when the
score is fixed before their labels exist. The participation-
volume and leakage audit
(subsection~\ref{sec:validation_leakage}) asks whether frequent
losers are overrepresented near defendants beyond what
participation volume mechanically predicts, and whether the
result survives stratification on participation history. The
direct-defendant asymmetry
(subsection~\ref{sec:validation_direct}) asks whether the
statistic is mistakenly behaving like a generic cartel-
membership detector; for a loser-side screen, it should not.

The section therefore evaluates the ranking against the threats
that would make it uninformative for enforcement: mechanical
exposure through participation volume, label leakage,
retrospective fitting, and legal overclaiming. The performance
criterion is not proof. It is whether award records rank
cartel-adjacent loser-side exposure well enough to allocate
costly bid-layer forensic effort.

\subsection{A Conservative Adjudication Benchmark}
\label{sec:validation_conservative}

The first benchmark is deliberately restrictive. We use only
CADE procurement-cartel cases adjudicated before the end of
2020, leaving $\valConservativeCases$ cases,
$\valConservativeFD$ firm-defendants active in BEC, and
$\valConservativeCobidders$ always-loser firms that bid
alongside at least one of those defendants. Of those cobidders,
$\valConservativeFL$ are frequent losers. Because both the
cartel finding and the defendant list precede the construction
of the screen, this benchmark anchors the initial validation
target in adjudications available early relative to the
temporal-holdout exercises that follow. The restriction trades
sample size for legal-timing discipline.

The main threat is mechanical exposure. Firms that participate
often have more opportunities to appear near any defendant,
regardless of cartel adjacency. Raw co-bidding rates would
therefore overstate what the screen contributes. The relevant
question is whether frequent losers are overrepresented near
adjudicated defendants beyond what their participation volume
alone would mechanically generate.

We address that concern with a participation-stratified
permutation. Within quartiles of participation volume, we
randomly permute cobidder status $\valPermPilot$ times and
recompute the frequent-loser share among placebo cobidders. The
procedure preserves the exposure created by participation
intensity while breaking the specific adjacency to CADE
defendants. The observed co-bidding rate is $\valCADEenrich$,
against a permutation mean of $\valCADEbase$, an excess ratio of
$\valMultThreeTwo\!\!\times$ with $p<0.001$. The ROC AUC against
the same conservative ground truth is $\valAUCprePost$ (95\%
DeLong CI $\valAUCprePostCI$).

This benchmark does not turn cobidder status into legal proof.
It answers the narrower validation question appropriate for a
loser-side screen: whether the ranking concentrates firms
adjacent to adjudicated cartel anchors at rates above what
participation volume mechanically predicts. The temporal-holdout
exercise in the next subsection extends the test to cases
adjudicated after the BEC sample closes, where the labels are
observed strictly after the score is constructed.

\subsection{Temporal Holdout Discrimination}
\label{sec:validation_temporal}

The broader CADE portfolio is more informative but less clean:
several linked cases are adjudicated after the BEC sample
closes. We treat it as a temporal holdout exercise rather than
as a pure contemporaneous benchmark. The labels are
later-adjudicated, so the exercise is not a real-time field
deployment. Its value is narrower: it asks whether a score
built from earlier award records ranks firms whose loser-side
cartel adjacency is revealed later by the adjudication record.

The holdout performance is strong. Computing the score from
$2009$--$2016$ participation and evaluating it on
$2017$--$2019$ labels gives firm-level AUC
$\valAUCFLfirmTemp$ (95\% CI $\valAUCFLfirmTempCI$). The
in-sample AUC is higher, $\valAUCFLfirm$, as expected. For an
enforcement screen that would be run before the case file is
complete, the holdout estimate is the more relevant benchmark.

\begin{figure}[htbp]
\centering
\includegraphics[width=0.82\textwidth]{./output/figures/fig_temporal_holdout_roc}
\caption[Temporal-holdout discrimination]{Temporal-holdout ROC
curves. Scores are computed using participation observed before
the test year; the test year supplies the out-of-time cobidder
label.}
\label{fig:temporal_holdout_roc_v18}
\end{figure}

Figure~\ref{fig:temporal_holdout_roc_v18} reports the
rolling-origin version of the same exercise: for each test
year, scores are constructed using only prior participation.
The relevant point is not that every year behaves identically,
but that the ranking remains separated from random
classification out of time.

The exercise supports temporal discrimination, not field
deployment, causal prediction, or legal proof. It addresses
retrospective fitting, since the score is fixed before the
test-period labels exist, but it does not eliminate every
timing concern: the labels are themselves observed later
through adjudication, and that informational content is
meaningful without being equivalent to real-time enforcement
evidence. The next subsection turns to a different threat:
whether the discrimination is driven by participation-volume
mechanics or by reuse of the same participation history.

\subsection{Participation Volume and Leakage}
\label{sec:validation_leakage}

The validation would be uninformative if it only reproduced
its own construction. Two mechanical concerns matter. The first
is \emph{exposure}: firms that participate often have more
opportunities to appear near any defendant, and the
participation-stratified permutation in
subsection~\ref{sec:validation_conservative} addresses that
margin. The second is \emph{leakage}: in item-level validation,
the same firms can help construct both the participation score
and the cobidder label. The two concerns are distinct and call
for different audits.

We therefore do not use the in-sample item-level AUC of
$\valLeakRefAUC$ as evidence of screening performance. That
number is a reference value for the mechanical upper bound: it
shows what happens before reuse of the same participation
history is disciplined. The credibility of the validation comes
from the audits that weaken that mechanical channel.

Two audits do the work. Holding cobidder firms out of score
construction within cross-validation folds leaves AUC
$\valLeakCVAUC$ (95\% CI $\valLeakCVAUCCI$); this addresses
reuse of the same firms in score and label construction.
Computing the score only on $2009$--$2016$ participation and
evaluating it against $2017$--$2019$ items leaves AUC
$\valAUCFLfirmTemp$ at the firm level; this addresses
retrospective fitting by placing the labels out of time.

The drop from the in-sample reference is real and disclosed; it
is the mechanical component. What remains is the relevant
signal for the screening interpretation: the participation
primitive continues to rank cartel-adjacent loser-side firms
when labels are held out of score construction or placed out of
time. The audits do not eliminate every source of mechanical
overrepresentation, but they isolate the surviving signal from
the channels they target. The next subsection turns to the
opposite check: whether the same statistic also classifies
direct CADE defendants, where a loser-side screen should
perform worse.

\subsection{Direct-Defendant Asymmetry as a Scope Check}
\label{sec:validation_direct}

The final validation check runs in the opposite direction. If
the award-layer statistic were a generic cartel-membership
detector, it should classify direct CADE defendants. A
loser-side ranking need not. The scope of the statistic predicts
asymmetric performance: it is built from zero wins and repeated
loser-side participation, while direct defendants are legal
membership anchors that often occupy winner-side or allocation
roles in the cartel environment.

That is what the data show. Against direct defendants in the
broader BEC firm universe, AUC is $\valAUCdirectStd$,
essentially random. We do not treat this as a hidden failure.
It is a scope check: the statistic is not behaving like a
general detector of cartel membership.

The interpretation is narrow. Direct defendants anchor
adjudicated cartel environments; cobidders are the loser-side
footprint the award layer can plausibly rank. Stronger
performance against cobidders than against direct defendants is
therefore not target shopping. It is the empirical signature of
the screen's stated scope.

The result rules out the wrong reading---that the screen
detects cartel membership generically---and supports the right
one: the screen ranks cartel-adjacent loser-side exposure
around adjudicated cartel environments. It should be read
together with the cobidder validation in the preceding
subsections, where the same statistic concentrates on the
population it was designed to rank.

\begin{table}[htbp]
\centering
\caption{Validation benchmarks for the loser-side ranking. Each
row addresses a distinct threat to interpretation.}
\label{tab:validation_benchmarks_v18}
\footnotesize
\begin{adjustbox}{max width=\textwidth,center}
\begin{tabular}{p{3.3cm}p{6.0cm}p{3.8cm}}
\toprule
Benchmark & Question & Result \\
\midrule
Conservative CADE cases & Are frequent losers overrepresented among pre-2020 adjudicated cobidders, net of participation volume? & $\valMultThreeTwo\!\times$ excess; AUC $\valAUCprePost$ \\
Temporal holdout & Does pre-window participation rank later cartel-adjacent cobidders? & AUC $\valAUCFLfirmTemp$ \\
Leakage audit & Does signal survive when cobidders are held out of score construction? & Out-of-fold AUC $\valLeakCVAUC$ \\
Direct-defendant scope check & Does the loser-side screen classify legal defendants symmetrically? & Direct-defendant AUC $\valAUCdirectStd$ \\
\bottomrule
\end{tabular}
\end{adjustbox}
\end{table}

Table~\ref{tab:validation_benchmarks_v18} gives the compact
view. The four rows are meant to be read together. The
conservative benchmark anchors the label in adjudicated cases;
the temporal holdout addresses retrospective fitting; the
leakage audit owns the mechanical component; and the
direct-defendant scope check keeps the legal object straight. Taken
together, they support a limited but useful claim: award records
can rank cartel-adjacent loser-side firms before bid-level proof
is recovered. They do not adjudicate membership, and they are not
asked to.
